\DAsection[Dockerfile 管单个服务，Compose 管整套应用]{Docker 部署 Node 前端与 FastAPI}

\begin{frame}[fragile]{8.1 静态前端 Dockerfile：构建阶段}{pnpm 安装依赖并生成 dist}
  \begin{lstlisting}[language=bash]
FROM node:24-slim AS build

ENV PNPM_HOME="/pnpm"
ENV PATH="$PNPM_HOME:$PATH"
RUN corepack enable

WORKDIR /app
COPY package.json pnpm-lock.yaml ./
RUN --mount=type=cache,id=pnpm,target=/pnpm/store \
    pnpm install --frozen-lockfile

COPY . .
RUN pnpm build
  \end{lstlisting}
  \DAtagline{缓存关键}{先复制 package 与 lockfile；源码变化不应让依赖安装层全部失效}
  \DAsource{pnpm Docs · Working with Docker}{https://pnpm.io/docker}
\end{frame}

\begin{frame}[fragile]{8.2 静态前端 Dockerfile：运行阶段}{运行镜像只包含 Nginx、配置与 dist}
  \begin{lstlisting}[language=bash]
FROM nginx:alpine AS runtime

COPY nginx.conf /etc/nginx/conf.d/default.conf
COPY --from=build /app/dist /usr/share/nginx/html

EXPOSE 80
CMD ["nginx", "-g", "daemon off;"]
  \end{lstlisting}
  \begin{itemize}
    \item 多阶段构建把开发依赖留在 build 阶段。
    \item 运行镜像只包含 Nginx、静态产物与配置，体积更小、攻击面更少。
    \item 前端的 Secret 不应在任何阶段注入；浏览器最终都能看到静态配置。
  \end{itemize}
\end{frame}

\begin{frame}[fragile]{8.3 前端容器内的 Nginx 配置}{负责 SPA 与容器网络内 /api 转发}
  \begin{lstlisting}[language=bash]
server {
    listen 80;
    root /usr/share/nginx/html;

    location / {
        try_files $uri $uri/ /index.html;
    }

    location /api/ {
        proxy_pass http://api:8000/;
        proxy_set_header Host $host;
        proxy_set_header X-Forwarded-Proto $scheme;
    }
}
  \end{lstlisting}
  \DAtagline{网络名}{\DAfile{api} 是 Compose 服务名，由 Docker DNS 在内部解析}
\end{frame}

\begin{frame}[fragile]{8.4 FastAPI Dockerfile：uv + 锁文件}{构建可重复的 Python 环境}
  \begin{lstlisting}[language=bash]
FROM python:3.12-slim

COPY --from=ghcr.io/astral-sh/uv:0.11.28 \
     /uv /uvx /bin/
WORKDIR /app

COPY pyproject.toml uv.lock ./
RUN uv sync --locked --no-dev --no-install-project

COPY . .
RUN uv sync --locked --no-dev

ENV PATH="/app/.venv/bin:$PATH"
CMD ["fastapi", "run", "app/main.py",
     "--host", "0.0.0.0", "--port", "8000",
     "--proxy-headers"]
  \end{lstlisting}
  \DAsource{uv Docs · Using uv with FastAPI}{https://docs.astral.sh/uv/guides/integration/fastapi/}
\end{frame}

\begin{frame}[fragile]{8.5 Compose 第一部分：数据库与 API}{变量从服务器 Env 文件注入}
  \begin{lstlisting}
services:
  db:
    image: mysql:8.4
    environment:
      MYSQL_DATABASE: ${MYSQL_DATABASE}
      MYSQL_USER: ${MYSQL_USER}
      MYSQL_PASSWORD: ${MYSQL_PASSWORD}
      MYSQL_ROOT_PASSWORD: ${MYSQL_ROOT_PASSWORD}
    volumes:
      - mysql_data:/var/lib/mysql
    healthcheck:
      test: ["CMD-SHELL",
             "mysqladmin ping -h 127.0.0.1 -uroot -p$$MYSQL_ROOT_PASSWORD"]
      interval: 10s
      timeout: 5s
      retries: 10
  \end{lstlisting}
\end{frame}

\begin{frame}[fragile]{8.6 Compose 第二部分：API 与前端}{前端映射到宿主机本地端口}
  \begin{lstlisting}[basicstyle=\ttfamily\fontsize{5.0}{5.6}\selectfont]
  api:
    build: ./backend
    env_file: .env.production
    depends_on:
      db:
        condition: service_healthy
    restart: unless-stopped

  frontend:
    build: ./frontend
    ports:
      - "127.0.0.1:8080:80"
    depends_on:
      - api
    restart: unless-stopped

volumes:
  mysql_data:
  \end{lstlisting}
  \DAtagline{暴露面}{API 与 MySQL 没有 \DAfile{ports}，只能从 Compose 网络访问}
\end{frame}

\begin{frame}[fragile]{8.7 第一次启动的顺序}{校验配置、构建、迁移、启动、验证}
  \begin{lstlisting}[language=bash]
docker compose config
docker compose build

docker compose up -d db
docker compose run --rm api \
  uv run alembic upgrade head

docker compose up -d
docker compose ps
docker compose logs --tail 100
  \end{lstlisting}
  \begin{itemize}
    \item 迁移工具以项目实际使用为准；核心是“只执行一次、失败就停止发布”。
    \item 首次启动后从宿主机测试 \DAcmd{curl http://127.0.0.1:8080/api/health}。
  \end{itemize}
\end{frame}

\begin{frame}{8.8 depends\_on 与健康检查}{启动顺序、健康状态与业务可用是三件事}
  \begin{itemize}
    \item 普通 \DAfile{depends\_on} 只控制容器启动顺序。
    \item 数据库容器启动后，MySQL 仍需要时间初始化。
    \item Healthcheck + \DAfile{condition: service\_healthy} 能减少早期连接失败。
    \item 应用自身仍应对短暂依赖故障实现有限重试与清晰日志。
  \end{itemize}
  \DAtagline{故障处理}{密码、迁移或网络错误应修复根因，避免无限重启}
\end{frame}

\begin{frame}[fragile]{8.9 宿主机 Nginx 接入 Compose}{宝塔继续负责域名与证书}
  \begin{lstlisting}[language=bash]
location / {
    proxy_pass http://127.0.0.1:8080;
    proxy_set_header Host $host;
    proxy_set_header X-Real-IP $remote_addr;
    proxy_set_header X-Forwarded-For $proxy_add_x_forwarded_for;
    proxy_set_header X-Forwarded-Proto $scheme;
}
  \end{lstlisting}
  \begin{itemize}
    \item 宝塔 Nginx 终止 HTTPS，再把整站请求转给前端容器。
    \item 前端容器内部 Nginx 再把 \DAfile{/api} 路由到 \DAfile{api:8000}。
    \item 两层代理便于课堂理解；生产也可把职责合并到一层 Nginx。
  \end{itemize}
\end{frame}

\begin{frame}[fragile]{8.10 Next.js / Node SSR}{运行 Node Server，由 Nginx 反向代理}
  \begin{lstlisting}[language=bash]
# next.config.ts
const nextConfig = { output: "standalone" };

# runtime stage
FROM node:24-slim AS runner
WORKDIR /app
ENV NODE_ENV=production
ENV HOSTNAME=0.0.0.0
ENV PORT=3000
COPY --from=builder /app/.next/standalone ./
COPY --from=builder /app/.next/static ./.next/static
COPY --from=builder /app/public ./public
CMD ["node", "server.js"]
  \end{lstlisting}
  \DAsource{Docker Docs · Containerize a Next.js application}{https://docs.docker.com/guides/nextjs/}
\end{frame}

\begin{frame}{8.11 Node 后端与 FastAPI 的 Compose 差异}{替换运行镜像与启动命令}
  \begin{DAtable*}{@{}p{25mm}p{45mm}p{50mm}@{}}
    \toprule
    项目 & FastAPI & Node.js API \\
    \midrule
    依赖锁 & \DAfile{uv.lock} & \DAfile{pnpm-lock.yaml} \\
    构建 & \DAcmd{uv sync --locked} & \DAcmd{pnpm install --frozen-lockfile} \\
    启动 & \DAcmd{fastapi run} & \DAcmd{node dist/main.js} / \DAcmd{pnpm start} \\
    容器端口 & 8000 & 3000 \\
    代理目标 & \DAfile{api:8000} & \DAfile{api:3000} \\
    \bottomrule
  \end{DAtable*}
  \DAtagline{架构复用}{域名、Nginx、数据库、Secret、日志、更新与回滚逻辑保持一致}
\end{frame}

\begin{frame}{8.12 pnpm Monorepo 部署}{只把目标应用与生产依赖带进镜像}
  \begin{itemize}
    \item 多包仓库可以在 build 阶段完成全量构建。
    \item 使用 \DAcmd{pnpm deploy --filter=<app> --prod <target>} 生成自包含部署目录。
    \item 每个服务从目标目录创建独立运行镜像。
    \item pnpm 11 默认要求 Workspace 正确配置 injected packages；升级前读版本文档。
  \end{itemize}
  \DAtagline{价值}{避免把整个 Monorepo、其他服务源码与全部开发依赖复制到每个镜像}
  \DAsource{pnpm Docs · pnpm deploy}{https://pnpm.io/cli/deploy}
\end{frame}

\begin{frame}[fragile]{8.13 Docker 更新流程}{重建镜像，原地替换容器}
  \begin{lstlisting}[language=bash]
git fetch origin
git switch main
git pull --ff-only

docker compose build
docker compose run --rm api \
  uv run alembic upgrade head
docker compose up -d

docker compose ps
docker compose logs --tail 100
  \end{lstlisting}
  \DAtagline{生产升级}{更成熟的流程由 CI 构建带 Commit SHA 的不可变镜像，服务器只负责 pull 与切换}
\end{frame}

\begin{frame}{8.14 Docker 回滚范围}{镜像、数据库、上传文件与配置}
  \begin{DAtable}{@{}p{27mm}p{89mm}@{}}
    \toprule
    对象 & 回滚策略 \\
    \midrule
    应用镜像 & 保留上一个稳定 Tag，Compose 切回旧 Tag \\
    数据库 Schema & 优先向后兼容迁移；破坏性迁移前做备份与演练 \\
    用户上传 & 独立 Volume / 对象存储，不随容器替换 \\
    环境变量 & 版本化记录变量名与变更，不提交真实值 \\
    Nginx / Compose & 配置进入 Git，变更前校验语法 \\
    \bottomrule
  \end{DAtable}
  \DAtagline{最高风险}{新代码写入了旧版本无法理解的数据时，单纯切回旧镜像可能仍然失败}
\end{frame}
